\scanpage{74}
% archive sheet number: 73
$A$ anneau
\[
\begin{aligned}
&\Delta\in A\\
&A_0=A/2A
\end{aligned}
\]
On veut écrire
\[
\Delta=b^2-4c\qquad b,c\in A \tag*{(1)}
\]
ce qui implique
\[
\Delta_0\in A_0^2\qquad(\Delta_0=(\Delta\bmod2))
\]
\deleted{soit}
\[
\unsure{\Delta_0|\omega_0}=b_0^2\qquad b_0\in A_0
\]
où $b_0$ est uniquement déterminé si \fbox{$A_0$ réduit}\textcircled{\scriptsize\unsure{a|*}} donc \deleted{\unsure{b}} $b$ déterminé mod $2$ \inserted{par la condition $b\equiv b_0\ (2)$}, donc si $b'$ satisfait la $=$ condition i.e $b'=b+2x$, on a
\[
{b'}^2=b^2+4(x^2+bx)\qquad\text{i.e.}\quad {b'}^{\unsure{2}}\equiv b^2\quad\text{\uuline{$(4)$}}
\]

On voit donc

\uline{Lemme} \quad $\forall\ b_0\in A_0$, il y a un unique élément \inserted{$\gamma(\deleted{\unsure{x}}b_0)$} de $A/4A$ tel que $\forall\ b\in A\deleted{\unsure{/2A}}$ avec $b\equiv b_0\ (2)$, on ait $\gamma(b_0)\equiv b^2\ (4)$. Pour $\Delta\in A$ donné, pour qu'il existe $b,c\in A$ tel qu'on ait (1), il f. et s. que

a) $\Delta_0\ (= \Delta\bmod2)$ soit le carré d'un él. $b_0$ de $A_0$ (bien déterminé car $A_0$ réduit)

b) $\gamma(b_0)=0$.

Supposons qu'on puisse trouver \deleted{un} $\unsure{v|\upsilon|u}\in A^{\deleted{\unsure{*}}}$ \deleted{\unsure{entier impair}} $\deleted{\unsure{u|v|\nu}=2\unsure{t|x|d}+1}$ \deleted{tel que les conditions précédentes sont}
